\documentclass[dvipsnames, svgnames, x11names]{beamer}

\usetheme{mytheme}
\usetikzlibrary{tikzmark}


%封面
\title{第二章\ 药物的吸收}
\author{张斌}
\institute[理工系]{河北农业大学 理工系}
\date{\the\year 年 \the\month 月}



\begin{document}

	
	\begin{frame}
		\titlepage
		\thispagestyle{empty}
		\addtocounter{framenumber}{-1}
	\end{frame}




% 药片口服以后，先在胃部崩解，药片中的化合物就会被释放出来；之后就开始随着胃的蠕动进入到肠。肠是主要的吸收器官，肠壁表面积很大，药物主要在这里被吸收。过肠的时候，肠壁细胞对药物会有代谢，过肠以后经肝门静脉进入肝脏，肝脏是最重要的代谢解毒器官，专门负责清除外源性分子，主要通过肝细胞中的代谢酶对药物分子进行各种各样的修饰，使其失活并被排出。幸存的未被肝脏萃取或修饰的药物分子才最终抵达心脏，经血液输送到全身。这是药物的吸收过程。

% 药物随血流遍布全身，会扩散进入几乎所有的器官和组织。这个过程称为分布。

% 药物在不断循环过程中，但凡经过肝脏，都会有一部分被肝脏细胞中的代谢酶修饰。这个过程称为代谢。

% 最终药物或者修饰过的药物经过胆汁排到肠中随粪便排出，或者经肾脏由尿排出，这是药物的排泄过程。

\begin{frame}
	\frametitle{ADME}
	\begin{minipage}{0.28\linewidth}
		药物的体内过程：
		\begin{itemize}
			\item Absorption ~吸收 
			\item Distribution 分布 
			\item Metabolism 代谢
			\item Excretion ~~~排泄
		  \end{itemize}
	\end{minipage}\hfill
	\begin{minipage}{0.68\linewidth}
		\begin{figure}[htpb]
			\centering
			\includegraphics[width=.9\linewidth]{ADME}
		\end{figure}
	\end{minipage}
\end{frame}


\begin{frame}{}
	\begin{center}
		{\Huge \cbf{目 \ 录}}
	\end{center}
	\tableofcontents[subsectionstyle=hide/hide/show]
	\thispagestyle{empty}
	\addtocounter{framenumber}{-1}
\end{frame}

\section{药物跨膜转运}

\begin{frame}
    \frametitle{药物跨膜转运}
    药物在体内吸收、分布、转化及排泄的过程中，首先必须跨越多层生物膜，进行多次转运，这种过程叫做药物跨膜转运。因此，了解药物跨膜转运机制及其影响因素是十分重要的。
\end{frame}

\subsection{生物膜的结构}

\begin{frame}
    \frametitle{经典模型}
    
    \begin{quote}
        生物膜=磷脂质+蛋白质+糖类
    \end{quote}

    1935年，Danielli与Davson提出细胞膜经典模型。
    
\begin{itemize}
    \item 细胞膜是由脂质双分子层构成；
    \item 两个脂质分子的疏水尾相连在中间形成膜的疏水区；
    \item 脂质分子的亲水头朝外分布在膜的外侧形成对称的双层膜结构；
    \item 膜蛋白分布在脂质双分子层的两侧；
    \item 膜上有许多带电荷的小孔，水分能自由通过；
    \item 膜结构中还存在特殊载体和酶促系统，能与某些物质特异性结合，进行物质转运。
  \end{itemize}
\end{frame}

\begin{frame}
    \frametitle{流动镶嵌模型}
    \hspace{-2em}
    \begin{minipage}{0.38\linewidth}

        1972年，Singer和Nicolson提出生物膜流动讓嵌模型，该模型强调了膜的流动性和不对称性，即膜的结构不是静止的而是流动的。

\begin{itemize}
  \item 该模型的基本结构仍是脂质双分子层
  \item 流动的脂质双分子层构成细胞膜的连续主体
  \item 蛋白质分子以不同的方式和不同的深度嵌入磷脂双分子层中
\end{itemize}

    \end{minipage}\hfill
    \begin{minipage}{0.6\linewidth}
        \begin{figure}[htpb]
            \centering
            \includegraphics[width=\linewidth]{流动镶嵌模型}
        \end{figure}
    \end{minipage}
    
\end{frame}

\begin{frame}
    \frametitle{晶格镶嵌模型}
    1975年由Wallach提出晶格镶嵌模型，进一步解释了膜的流动性和完整性特征
    \begin{itemize}
      \item 膜的流动性是由于脂质能可逆地进行无序（液态）和有序（晶态）的相变过程
      \item 膜蛋白对脂质分子的活动具有控制作用
      \item 具有流动性的脂质是呈小片的点状分布
      \item 脂质的流动性是局部的，并不是整个脂质双分子层都在流动
    \end{itemize}
    解释了为什么细胞膜既具有流动性又能保持其完整性和稳定性。
\end{frame}


\begin{frame}
    \frametitle{生物膜的性质}
    \begin{itemize}
      \item 流动性
      \item 不对称性
      
      细胞膜内外两侧层面的组分和功能有明显的差异
      \item 选择透过性
      
      细胞膜具有选择透过性 ，可以让水分子自由透过 ，选择吸收的离子和小分子也可以通过 ，而其他的离子 、小分子和大分子则不能通过 。例如 ，膜的脂质结构特征 ，使得脂溶性药物较易透过而脂溶性很小的药物难以透过 ；膜上特异性表达的载体蛋白是某些药物选择性透过的载体 。
    \end{itemize}
\end{frame}


\subsection{药物的跨膜转运方式}

\begin{frame}
    \frametitle{药物的跨膜转运方式}
\begin{figure}[htpb]
    \centering
    \includegraphics[width=.78\linewidth]{自由扩散}
\end{figure}
\end{frame}

\begin{frame}
    \frametitle{被动扩散(单纯扩散)}
    被动转运是指不需要消耗能量，生物膜两侧的药物由高浓度侧向低浓度侧转运的过程。大多数药物是通过这种方式转运的，即药物是借助于在生物膜中的脂溶性顺浓度差跨膜转运的。
    
        \begin{enumerate}
          \item 顺浓度梯度转运，可以用Fick定律描述药物的转运速率
          \[
            \frac{\mathrm{d}C}{\mathrm{d}t}=\frac{DAk}{h}(C_{\mathrm{GI}}-C)
          \]
          $\frac{\mathrm{d}C}{\mathrm{d}t}$为扩散速率；$D$ 为扩散系数；$A$为扩散表面积， $k$为油水分配系数；$h$为膜厚度 ；$C_{\mathrm{GI}}$为胃肠道中的药物浓度 ；$C$为血药浓度。
          \item 膜对转运的药物无选择性，转运速率只与其脂溶性大小及扩散系数有关
          \item 膜对转运的药物量无饱和现象
          \item 无竞争性抑制作用
        \end{enumerate}
\end{frame}


\begin{frame}
    \frametitle{被动扩散(单纯扩散)}
    \[
        \frac{\mathrm{d}C}{\mathrm{d}t}=\frac{DAk}{h}(C_{\mathrm{GI}}-C)
    \]
    当药物口服后 ，胃肠道中的浓度远大于血中的药物浓度，则$C$可以忽略不计；在给予某一药物于某一个体的吸收过程中，其 $D$、$A$、$h$、$k$都为定值 ，可用渗透系数$P$ (permeability coefficient) 来表达 ， 即$ P =DAk/h$ ， 此时
    \[
        \frac{\mathrm{d}C}{\mathrm{d}t}=PC_{\mathrm{GI}}
    \]
    即药物的扩散速率等于渗透系数与胃肠道药物浓度的乘积 。渗透系数是药物渗透性的参数，描述药物的膜渗透能力。对于给定药物，药物单纯扩散透过膜的转运速率与胃肠道中药物浓度呈线性关系 。
\end{frame}

\begin{frame}
  \frametitle{促进扩散}
  又称易化扩散、协助扩散，或帮助扩散。

  促进扩散是指某些药物在细胞膜上转运体的帮助下 ， 由高浓度侧向低浓度侧跨细胞膜转运的过程 。有些药物虽然水溶性不好 、 脂溶性也较差 ，但也能有较好的透膜吸收， 这是因为膜结构中的一些特殊转运体参与了药物转运 。 一般认为 ， 促进扩散的转运机制是 ： 细胞膜上的转运体在膜外侧与药物结合后 ， 通过转运体的自动旋转或变构将药物转运到细胞膜内侧 。

  \begin{enumerate}
    \item 速率快 、 效率高
    \item 有选择性
    \item 有饱和现象
    \item 有部位特异性
    \item 有竞争性抑制现象
  \end{enumerate}

\end{frame}


\begin{frame}
    \frametitle{孔道转运}
    生物膜上有水通道或蛋白质分子孔道。在毛细血管壁上也存在细胞间隙构成的孔道。一些物质包括水和某些电解质等可以通过这些孔道转运。这种转运通常与药物的分子结构和大小有关。转运速率主要取决于相应组织的血流速率以及生物膜的性质，而与脂溶性和pH的梯度关系不大。但脑等特殊组织，由于毛细管内皮细胞紧密联接，缺乏孔道转运。 

    \begin{enumerate}
      \item 主要为水和电解质的转运
      \item 转运速率与所处组织及膜的性质有关
    \end{enumerate}
\end{frame}


\begin{frame}
    \frametitle{主动转运}
    \begin{figure}[htpb]
      \centering
      \includegraphics[width=.9\linewidth]{主动运输}
    \end{figure}
\end{frame}


\begin{frame}
    \frametitle{主动转运}
    主动转运是指生物膜两侧的药物借助载体蛋白的帮助由低浓度侧向高浓度侧（ 逆浓度梯度 ） 转运的过程 。主动转运需要消耗能量。

    主动转运速率可用米氏方程描述：
    \[
        \frac{\mathrm{d}C}{\mathrm{d}t}=\frac{V_mC}{K_m+C}
    \]
    主动转运可分为 ATP 驱动泵和协同转运两种 。
    \begin{enumerate}
      \item ATP 驱动泵
      
      以 ATP 水解释放的能量为直接能源进行主动转运的载体蛋白家族称为ATP 驱动泵 。 这类载体蛋白也是一种 ATP 酶 ，能催化 ATP 水解提供能量 。
      \item 协同转运
      
      转运器不直接利用 ATP 水解的能量 ， 而是借助膜上相邻钠钾泵排出Na$^+$ 所产生的势能贮备，与 Na$^+$ 相伴或后继进行药物转运的主动转运方式。
    \end{enumerate}
\end{frame}


\begin{frame}
    \frametitle{主动转运特点}
    \begin{enumerate}
      \item 逆浓度梯度转运
      \item 常需要能量
      \item 需要载体参与
      \item 转运速率有饱和性
      \item 有竞争性抑制现象，同类物质往往能竞争同一载体, 产生相互拮抗作用。
      \item 对物质转运存在化学和立体选择性
    \end{enumerate}
\end{frame}


\begin{frame}
    \frametitle{膜动转运}
    转运蛋白虽然能够帮助许多离子和小的分子通过细胞膜，但是，对于像蛋白质和多糖这样的生物大分子的运输却无能为力。
    \begin{figure}[htpb]
      \centering
      \includegraphics[width=.6\linewidth]{胞吞与胞吐}
    \end{figure}
\end{frame}


\begin{frame}
    \frametitle{膜动转运}
    通过细胞膜的主动变形将物质摄入细胞内或从细胞内释放到细胞外的转运过程。
    \begin{enumerate}
      \item 入胞作用
      
      物质借助与细胞膜上某些蛋白质的特殊亲和力而附着于细胞膜上，通过细胞膜的内陷形成小泡，包裹药物的小泡逐渐与细胞膜表面断离，从而将物质摄入细胞内的转运过程

      \begin{itemize}
        \item 蛋白质 、 多肽类 、 脂溶性维生素
        \item 重金属
      \end{itemize}
      \item 出胞作用
      
      某些大分子物质通过形成小泡从细胞内部移至细胞膜内表面，小泡的膜与细胞膜融合 ，从而将物质排出细胞外的转运过程

      \begin{itemize}
        \item 腺细胞分泌胰岛素
        \item 细胞内不能消化的物质
      \end{itemize}
    \end{enumerate}
\end{frame}

\begin{frame}
    \frametitle{药物的转运机制是一个非常复杂的过程}
    药物以何种机制转运吸收 ， 与药物性质 、 吸收部位生理特征等密切相关 。 
    
    某种药物可以通过一种特定的转运机制吸收 ， 也可以通过多种转运机制吸收 。
\end{frame}


\begin{frame}[t]
    \frametitle{}
    \begin{question}
        药物借助载体或酶促系统，消耗机体能量，从膜的低浓度向高浓度一侧转运的药物转运方式是
        \tcblower

        \only<1>{
            A. 主动转运
        }
        \only<2>{
            \highlight{A. 主动转运}
        }

        B. 简单扩散

		    C. 易化扩散

        D. 膜动转运

        
    \end{question}
\end{frame}


\begin{frame}[t]
    \frametitle{}
    \begin{question}
        扩散速度取决于膜两侧药物的浓度梯度，药物的脂水分配系数及药物在膜内扩散速度的药物转运方式是
        \tcblower

        A. 滤过

        \only<1>{
            B. 简单扩散
        }
        \only<2>{
            \highlight{B. 简单扩散}
        }

        C. 易化扩散
        
        D. 主动转运

        
    \end{question}
\end{frame}


\section{胃肠道的结构与功能}

\begin{frame}
	\frametitle{药物的吸收}

	\highlight{吸收}是指药物从给药部位进入血液循环的过程。除了动脉和静脉给药物外，其它给药途径均存在吸收过程。

	\highlight{生物利用度} 药物从给药部位进入血液循环过程通常用吸收速度和吸收程度来描述。药物吸收程度通常指生物利用度（bioavailability），即药物由给药部位到达血液循环中的相对量。

	由于口服给药的部位是胃、小肠和大肠，故称为\highlight{胃肠道吸收}。口服给药方式在临床应用最为广泛，因口服给药安全、方便，患者的顺应性较好。了解胃肠道的功能及与吸收有关的生理特征，有利于掌握口服药物吸收的规律。

\end{frame}


\begin{frame}
    \frametitle{胃肠道的结构}
    \hspace{-2em}
    \begin{minipage}{0.28\linewidth}
        \begin{figure}[htpb]
            \centering
            \includegraphics[width=\linewidth]{胃肠道}
        \end{figure}
    \end{minipage}\hfill
    \begin{minipage}{0.68\linewidth}
        \begin{table}
            \centering
            \begin{tabular}{ccccc}
                \toprule
               部位 & pH & 长度(cm) & 表面积 & 转运时间 \\ \midrule
               胃 & $1 \sim 4$ & --- & 小 & $0.5 \sim 3$ h\\
               十二指肠 & $4 \sim 6$ & $20 \sim 30$ & 较大 & \multirow{3}{*}{$3 \sim 4$ h}\\
               空肠 & $6 \sim 7$ & $150 \sim 250$ & 很大 & \\
               回肠 & $6.5 \sim 7.5$ & $200 \sim 350$ & 很大 & \\
               盲肠/右结肠 & $5.5 \sim 7.5$ & \multirow{2}{*}{$90 \sim 150$} & \multirow{2}{*}{较小} & \multirow{2}{*}{$14 \sim 80$ h}\\
               左结肠/直肠 & $6.4 \sim 7.5$ &  &  & \\ \bottomrule
            \end{tabular}
        \end{table}

        因消化道各部位组织结构以及相应的pH值不同，对药物的吸收能力与吸收速度也是不同的。
    \end{minipage}
\end{frame}


\begin{frame}
    \frametitle{胃}
    \begin{itemize}
      \item 胃是消化道中最膨大的部分，成人的容量一般约为$1 \sim 2$升
      \item 胃液的pH值约为$0.9 \sim 1.5$ \footnote{正常胃液的pH为$0.9 \sim 1.5$，饭后胃液被稀释，pH值上升}
      \item 胃液呈强酸性，除一些弱酸性药物有较好吸收外，大多数药物胃内吸收较差
      \item 胃的吸收表面积小，血流速率也小，加之在胃中停留的时间较短，因此，胃不是药物的主要吸收部位
      \item 某些酸不稳定的药物可能因分解而失活。在这种情况下，应考虑将药物制成肠
      溶片，以避免胃酸对药物的分解作用。  
    \end{itemize}
\end{frame}


\begin{frame}
    \frametitle{小肠}
    \begin{itemize}
      \item 小肠由十二指肠 、 空肠和回肠组成。人的小肠长约4米。
      \item 小肠液的 pH 为$5 \sim 7.5$， 是弱碱性药物吸收的最佳环境。
      \item 小肠粘膜具有环形皱褶，并有大量的绒毛。绒毛外面是一层柱状上皮细胞，每个柱状上皮细胞的顶端约有 1700条微绒毛。
      \item 由于环形皱褶、绒毛和微绒毛的存在，最终使小肠的吸收面积比同样长的简单圆筒面积增加约600倍，达200 m$^2$左右。因而小肠是药物吸收过程进行的区域。
      \item 一般来说，糖、氨基酸和脂肪以及大部分药物是在十二指肠和空肠吸收的。
      \item 物质在小肠内的转运时间约$3 \sim 4$ h
    \end{itemize}
\end{frame}


\begin{frame}
    \frametitle{大肠}
    \begin{itemize}
      \item 大肠是由盲肠 、 结肠 （ 升结肠 、 横结肠 、 降结肠 、 乙状结肠 ） 和直肠组成 。
      \item 大肠的主要功能是储存食物糟粕 、 吸收水分 、 无机盐及形成粪便 。
      \item 大肠粗而短而粗， 黏膜上有皱纹但无绒毛 ，因而有效吸收表面积比小肠小得多 ， 药物吸收也比小肠差 。
      \item 除结肠定位给药和直肠给药外 ，只有一些吸收很慢的药物 ， 在通过胃与小肠未被吸收时 ， 才呈现药物吸收功能 。
      \item 物质在小肠内的转运时间约$3 \sim 4$ h，在大肠中需要$10 \sim 12$ h 。多数药物往往存在最佳吸收部位 ， 即吸收窗。当药物在胃肠中的释放时间大于 4 h或更长 ， 肯定会有一部分药物在大肠中释放 。
    \end{itemize}
\end{frame}


\begin{frame}
    \frametitle{药物胃肠转运机制}
    \hspace{-2em}
    \begin{minipage}{0.3\linewidth}
        \begin{enumerate}
            \item 药物通过不搅动水层
            \item 药物通过肠上皮
            \item 药物透过细胞间隙
            \item 通过淋巴吸收
          \end{enumerate}
    \end{minipage}\hfill
    \begin{minipage}{0.68\linewidth}
        \begin{figure}[htpb]
            \centering
            \includegraphics[width=\linewidth]{胃肠转运机制}
        \end{figure}
    \end{minipage}
\end{frame}


\section{影响药物吸收的因素}

\begin{frame}
	\frametitle{影响药物吸收的因素}
    \begin{enumerate}
      
      \item 生理因素
        \begin{itemize}
          \item 胃肠pH、 胃肠活动性
          \item 肝功能及肝肠血流灌注情况
          \item 胃肠结构和肠道菌丛状况
          \item 年龄、性别、遗传因素
          \item 病人饮食特点
        \end{itemize}
    \item 药物因素
        \begin{itemize}
            \item 药物理化性质（如粒径大小、溶解度和药物的晶型等）
        \end{itemize}
    \item 制剂因素
    \begin{itemize}
        \item 处方中赋型剂的性质与种类、制备工艺、药物的剂型等
    \end{itemize}
    \end{enumerate}
\end{frame}


\subsection{生理因素}


\begin{frame}
    \frametitle{胃肠液的成分与性质}
    \begin{itemize}
      \item 胃液的主要成分是胃酸 （ 盐酸 ） ， 正常成人每日分泌的胃液量为1.5 $\sim$ 2.5L ， 空腹时胃液 pH 为 $0.9\sim 1.5$  ， 饮水或进食后 ，pH 可上升至$3.0\sim 5.0$ 。 由于胃液的 pH呈酸性 ， 有利于弱酸性药物的吸收 ， 而弱碱性药物吸收甚少 。 胃液 pH 的改变对弱酸性药物的吸收有影响 。
      \item 小肠自身分泌液是一种弱碱性液体 ， pH 约为 7.6 ， 成人每天分泌量1 $\sim$ 3L 。 小肠较高的 pH 环境是弱碱性药物最佳的吸收部位 。 
      \item 胃肠道中酸 、 碱性环境可能对某些药物的稳定性产生影响 。
      \item 胃肠液中含有酶类 、 胆盐等物质 ， 对药物的吸收产生不同的影响 。 胃蛋白酶 、 胰酶等可以消化食物 ， 也能分解多肽及蛋白质 ， 因此多肽与蛋白质药物口服易分解而失效 。
      \item  胆汁中含有胆酸盐 ， 是一种表面活性剂 ，能增加难溶性药物的溶解度 ，从而提高这类药物的吸收速度和程度 ； 胆酸盐也能与一些药物形成难溶性盐 ， 从而降低药物吸收
    \end{itemize}
\end{frame}


\begin{frame}
    \frametitle{胃排空和胃空速率}
    胃既有储存食物的功能，又具有“泵”的作用。口服给予的药物由胃至小肠的转运速度或时间，对药物吸收的速度和程度有显著影响。

    \begin{enumerate}
      \item 胃排空 ： 胃内容物从胃幽门排入十二指肠的过程称为胃排空。
      \item 胃空速率 ： 胃排空的快慢用胃空速率来描述 。
      
      胃排空按照一级速率过程进行 ， 可用胃空速率常数或胃空半衰期来表达 ， 服从下式\footnote{$\mathrm{ln(x)=2.303}\ \mathrm{lg(x)}$} ：
      \[
        \mathrm{lg}V_{t}=\mathrm{lg}V_{0}-\frac{K_{\mathrm{em}}}{2.303}t
      \]
      式中，$V_t$为$t$时间胃内容物体积 ； $V_{0}$为初始时胃内容物体积 ；$K_{\mathrm{em}}$为胃空速率常数 。
    \end{enumerate}
    
    
\end{frame}


\begin{frame}
    \frametitle{胃排空和胃空速率}
    胃排空的快慢对药物在消化道中的吸收有一定影响 。 
    \begin{itemize}
      \item 胃排空速率慢 ， 药物在胃中停留时间延长 ， 与胃黏膜接触机会和面积增大 ， 主要在胃中吸收的弱酸性药物吸收会增加 。 
      \item 由于大多数药物的主要吸收部位在小肠 ，因此 ，胃空速率决定了药物到达肠道的速度 ，对药物的起效快慢 ， 药效强弱及持续时间有显著的影响 。 
      \item 少数在特定部位吸收的药物 ， 胃空速率大时吸收反而较差
      \item 对于一些会被胃酸或酶降解的药物 ，胃排空迟缓将增加药物的降解程度。
    \end{itemize}
\end{frame}


\begin{frame}
    \frametitle{影响胃空速率的因素}
    \begin{figure}[htpb]
        \centering
        \includegraphics[width=.75\linewidth]{因素}
    \end{figure}
\end{frame}


\begin{frame}
    \frametitle{食物的影响}
    \begin{enumerate}
      \item 延缓或减少药物的吸收 ：
      
      食物除了改变胃空速率而影响吸收外 ， 食物还能消耗胃肠内水分 ， 使胃肠黏液减少 ， 固体制剂的崩解 、 药物的溶出变慢 ， 从而延缓药物的吸收 。 食物的存在还可增加胃肠道内容物的黏度 ， 使药物的扩散速度减慢而影响吸收 。 
      \item 促进药物的吸收 ：
      
      脂肪类食物具有促进胆汁分泌的作用 ， 由于胆汁中的胆酸离子具有表面活性作用 ， 可增加难溶性药物的溶解度而促进其吸收 。 有部位特异性吸收的药物可因食物降低胃空速率而增加吸收。 
    \end{enumerate}
\end{frame}


\begin{frame}
    \frametitle{食物的影响}
    \begin{figure}[htpb]
        \centering
        \caption{进食对药物吸收的影响}
        \includegraphics[width=.9\linewidth]{食物}
    \end{figure}
\end{frame}


\begin{frame}
    \frametitle{肝首过效应}
    \highlight{首过效应}(first pass effect)口服给药，药物在尚未吸收进入血循环之前，经肠道、肠壁和肝脏的代谢分解，使进入体内的相对药量降低的现象。

    \medskip

    胃吸收~~~~ → |~~~~~~~~~~~~ 胃冠状静脉、胃网膜左静脉 ~~~~~~~~~~~|  →

    小肠吸收 → | 十二指肠静脉、小肠静脉、肠系膜上静脉|  → 肝门静脉 → 肝脏 → 体循环

    大肠吸收 → |~~~~~~~~~ 肠系膜上静脉、肠系膜下静脉~~~~~~~~~~~|  →

    \medskip

    口服给药必须经胃肠道（壁）和肝脏后才进入体循环，有些药物几乎无代谢作用发生，有些则在胃肠壁或肝脏内被广泛代谢、消除，发生首过作用。
    
    对于首过效应大的药物，口服给药往往生物利用度很低， 或个体差异大，难以获得满意的疗效。对于这类药物最好采取其它给药途径，避免首过效应。 
\end{frame}


\begin{frame}
    \frametitle{疾病}
    \begin{enumerate}
      \item 疾病引起的胃肠道 pH 改变能影响药物从剂型中的溶出及吸收 。
        \begin{itemize}
          \item 胃癌患者的胃液 pH 往往升高
          \item 酸分泌长期减少的贫血患者服用铁剂及西咪替丁时吸收缓慢
        \end{itemize}
      \item 肝脏疾病常伴有其他脏器功能的变化 ， 从而造成对药物体内过程的影响 。
        \begin{itemize}
          \item 肝硬化患者由于肝细胞活性下降及合并门静脉旁路 ， 使相当多的胃肠道血液绕过门脉循环而通过门脉外循环直接进入体循环 ， 这样可引起某些首过代谢程度高的药物的口服生物利用度增加 。
        \end{itemize}
      \item 甲状腺功能异常可能影响药物的吸收
        \begin{itemize}
          \item 甲状腺功能不足时，肠的转运速率往往降低，使维生素 B2 在小肠吸收部位滞留的时间延长
          \item 甲状腺功能亢进的儿童则对维生素 B2 的吸收减少
        \end{itemize}
    \end{enumerate}
\end{frame}

\subsection{药物因素}

\begin{frame}
    \frametitle{药物的解离度}
    在胃肠道液中已溶解的弱酸或弱碱性药物以未解离型 （ 分子型 ） 和解离型两种形式存在 ， 两者所占比例由药物的解离常数 $\mathrm{p}K_\mathrm{a}$ 、 和胃肠道吸收部位 pH 所决定 。由于胃肠道上皮细胞膜为类脂膜 ， 通常脂溶性较大的未解离型药物容易通过被动扩散机制透过上皮细胞膜被吸收 ， 而解离后的离子型药物不易透过 ， 难以吸收 。

    胃肠液中未解离型与解离型药物浓度之比可用 Henderson-Hasselbalch 方程式描述 ：
    \[
        \text{弱酸性药物}:\mathrm{p}K_\mathrm{a}\mathrm{-p}\mathrm{H}=\mathrm{lg~}\frac{C_\mathrm{u}}{C_\mathrm{i}}
    \]
    \[
        \text{弱碱性药物}:\mathrm{p}K_\mathrm{a}\mathrm{-p}\mathrm{H}=\mathrm{lg~}\frac{C_\mathrm{i}}{C_\mathrm{u}}
    \]
    式中$C_\mathrm{u}$， $C_\mathrm{i}$分别为未解离型和解离型药物的浓度 。

\end{frame}


\begin{frame}[t]
    \frametitle{}
    \begin{example}
        弱酸性药物水杨酸的 $\mathrm{p}K_\mathrm{a}$ 为 3.0 ， 在 pH 为 1 的胃液中未解离型与解离型的比例为？在pH 为 7 的肠液中该比例为？        
    \end{example}

    \pause
    \qquad (1) 100：1

    \qquad (2) 1：10000
    
    \pause
    \begin{example}
        弱碱性药物奎宁的 $\mathrm{p}K_\mathrm{a}=8.4$， 在 pH 为 7 的肠液中未解离型与解离型的比例为1 ： 25， 在 pH 为 1 的胃液中该比例为 1 ：(\num{2.5 e7}) 。
    \end{example}

    通常弱酸性药物在胃中有较好的吸收 ， 而弱碱性药物在肠中有较好的吸收 。
\end{frame}


\begin{frame}
    \frametitle{药物的解离度}
    \begin{figure}[htpb]
        \centering
        \includegraphics[width=\linewidth]{非离子型}
        \caption{几种酸性 (A) 和碱性 (B) 药物在不同肠腔中的非离子型百分数比较}
    \end{figure}
\end{frame}


\begin{frame}
    \frametitle{药物的脂溶性}
    胃肠道黏膜上皮细胞为类脂膜 ， 而细胞外是水性环境 ， 因此药物分子若要通过被动扩散渗透进入细胞 ， 必须具有合适的水溶性和脂溶性 。 评价药物脂溶性大小的参数是油 / 水分配系数($K_\mathrm{o/w}$) 。 药物穿越细胞的能力与它的油 / 水分配系数存在相关性。 

    \hspace{-2em}
    \begin{minipage}{0.48\linewidth}
        \begin{itemize}
            \item 通常药物的$K_\mathrm{o/w}$大 ， 说明该药物的脂溶性较好 ， 吸收率也大
            \item 被动扩散机制吸收还与分子量相关 ， 分子量较小的药物更易穿透生物膜
          \end{itemize}
    \end{minipage}\hfill
    \begin{minipage}{0.48\linewidth}
        \begin{figure}[htpb]
            \centering
            \includegraphics[width=\linewidth]{吸收率}
        \end{figure}
    \end{minipage}
    
\end{frame}


\begin{frame}
    \frametitle{药物的溶出}
    \begin{minipage}{0.48\linewidth}
        对于多数固体剂型来说，药物的溶出速度直接影响药物的吸收速度。药物的溶解过程包括两个连续的阶段，首先是溶质分子从固体表面溶解形成饱和层，然后在扩散作用下经过扩散层，再在对流作用下进入溶液主体内。
    \end{minipage}\hfill
    \begin{minipage}{0.38\linewidth}
        \begin{figure}[htpb]
            \centering
            \includegraphics[width=\linewidth]{溶出}
        \end{figure}
    \end{minipage}
    
    固体药物的溶解速度主要受扩散控制，可用Noyes-Whitney 方程表示:    
    \[
        {\frac{\mathrm{d}C}{\mathrm{d}t}}={\frac{DS}{Vh}}(C_{\mathrm{s}}-C)
    \]

    式中，$\frac{\mathrm{d}C}{\mathrm{d}t}$为溶出速度；D为溶质在溶出介质中的扩散系数；$S$为固体的比表面积；$V$为溶出介质的体积；$h$为扩散层的厚度；$C_\mathrm{s}$为溶质在溶出介质中的溶解度；$C$为$t$时间溶液中溶质的浓度。
\end{frame}


\begin{frame}
    \frametitle{影响药物溶解速度的因素}
    \begin{enumerate}
      \item 固体的表面积
      
      固体药物的粒径越小，孔隙率越大，表面积越大，药物的溶解速度越快。
      \item 扩散系数
      
      在温度一定的条件下，扩散系数大小受溶出介质的黏度和药物分子大小的影响。
      \item 扩散层的厚度
      
      可通过搅拌，减小扩散层的厚度，加快药物溶出速度。
      \item 溶出介质的体积
      
      溶出介质的体积小，溶液中药物浓度高，药物溶出速度慢；反之，则药物溶出速度快。
      \item 温度
      
      温度升高，药物溶解度$C_\mathrm{s}$增大，扩散增强，黏度降低，药物溶出速度加快。
    \end{enumerate}
\end{frame}


\begin{frame}
    \frametitle{}
    \begin{question}
        固体药物的粒径越小，药物的溶解速度越快？还是粒径越大溶解速度越快？
    \end{question}\pause
    \begin{figure}[htpb]
        \centering
        \includegraphics[width=.48\linewidth]{粒径}
        \caption{非那西丁颗粒粒径与血药浓度的关系}
    \end{figure}
\end{frame}



\begin{frame}
    \frametitle{药物在胃肠道中的稳定性}
    胃肠道分泌液 、 不同 pH 、 消化酶 、 肠道菌群及细胞内代谢酶等 ， 可使口服药物在吸收前产生
降解或失去活性 ， 因而在药物剂型设计 、 制剂处方工艺设计时应加以注意 。

\begin{enumerate}
  \item 注射
  \item 为提高药物在胃肠道中的稳定性 ， 可制成药物的衍生物或前体药物。
  \item 采用肠溶包衣技术防止药物在胃酸中降解 。
\end{enumerate}
\begin{blankblock}{}
    青霉素极易水解 ， 在 pH 为 1 时半衰期仅为 33 秒 ， 故不宜口服而设计成注射用粉针 ；

    胰岛素极易被胃肠道消化酶破坏失去活性 ， 且分子量大不易被吸收 ， 故设计成注射剂 。

    青霉素的衍生物氨苄西林 ， 在 pH 为 1 时的半衰期为 5 小时 ， 即使在胃酸中也远比青霉素稳定 ， 因而可设计成口服制剂。
\end{blankblock}

\end{frame}


\subsection{剂型与制剂因素}

\begin{frame}
    \frametitle{\thesection.\thesubsection.1 剂型与药物吸收}
    药物的剂型对药物的吸收及生物利用度有很大的影响 。 不同药物的剂型 ， 给药部位及吸收
途径各异 ， 药物吸收的速度与程度亦可能不同 。

    剂型中药物的吸收和生物利用度情况取决于剂型释放药物的速度与数量 。 一般认为 ， 口服
剂型生物利用度高低的顺序为 ： 溶液剂 > 混悬剂 > 颗粒剂 > 胶囊剂 > 片剂 > 包衣片 。

\begin{figure}[htpb]
    \centering
    \includegraphics[width=.48\linewidth]{剂型因素}
\end{figure}
\end{frame}


\begin{frame}
    \frametitle{液体剂型}
    \begin{enumerate}
      \item 溶液型药物以分子或离子状态分散在介质中 ， 因而口服溶液剂的吸收是口服剂型中最快 、 且较完全的 ， 生物利用度高 。
      \item 乳剂的口服生物利用度较高 。
      \item 混悬剂中药物颗粒必须溶解才能被吸收 ， 溶解过程是否为吸收的限速过程取决于药物的溶解度和溶出速率 。
    \end{enumerate}
\end{frame}


\begin{frame}
    \frametitle{固体剂型}
    \begin{minipage}{0.58\linewidth}
        \begin{enumerate}
            \item 散剂比表面积大 ， 易分散 ， 服用后可不经崩解和分散过程 ， 所以吸收较其他固体口服制剂快 ， 生物利用度较高 。 
            \item 胶囊剂制备时一般填充颗粒或粉末 ， 因而囊壳崩解后药物内容物可快速分散 、 溶出 ， 故药物吸收快 ， 吸收较好 。
            \item 片剂是广泛应用的剂型之一 。 片剂在胃肠道中经历着崩解 、 分散和溶出的全过程 。 片剂充分崩解 ， 分散成包含辅料的细颗粒 ， 细颗粒进一步分散 ， 药物溶解后才能被机体吸收。
          \end{enumerate}
    \end{minipage}\hfill
    \begin{minipage}{0.33\linewidth}
        \begin{figure}[htpb]
            \centering
            \includegraphics[width=\linewidth]{片剂}
        \end{figure}
    \end{minipage}
\end{frame}


\begin{frame}
    \frametitle{\thesection.\thesubsection.2 辅料的影响}
    制剂的处方组成 ， 包括药物来源及剂量 、 辅料种类和加入量 ， 对药物的吸收均有影响 。 不同厂家制备的同一种药物制剂 ， 由于处方组成不同 ， 制剂的体外质量和口服生物利用度可有较大的差异 。
    
    制剂中常需要添加各种辅料 ， 以获得满意的加工特性、 良好的稳定性和期望的制剂学特性。 无活性的辅料几乎不存在 ， 许多辅料对药物和制剂可产生多方面影响 。
    \begin{itemize}
      \item 黏合剂
      \item 稀释剂
      \item 崩解剂
      \item 润滑剂
      \item 表面活性剂
    \end{itemize}

\end{frame}


\begin{frame}
    \frametitle{稀释剂}
    片剂在压片时要求物料有一定的数量和具有一定的性质才能成型，所以需要加入赋形剂。特别是小剂量的药物，为方便制剂操作，避免制剂过程中药物的损耗，必须加入稀释剂。

    药物与稀释剂之间常见的相互作用主要是稀释剂对主药的吸附作用或分散作用 。 若稀释剂为不溶性物质而又有较强的吸附作用 ，则被吸附的药物很难释放出来 ， 其生物利用度会显著降低。
    \begin{blankblock}{}
        三硅酸镁和碳酸镁能吸附抗胆碱药物阿托品等 。 
        
        疏水性药物中加入亲水性分散剂可对药物起到较好的分散作用 ， 能够减少药物粉末与液体接触时的结块现象 ， 使药物有合适的有效比表面积 ， 有利于药物的吸收 。
    \end{blankblock} 
\end{frame}


\begin{frame}
    \frametitle{黏合剂}
    片剂制粒过程中常常加入黏合剂以增加颗粒之间的黏结能力 ， 便于制粒 ， 但过量使用可延缓片剂的崩解 。 选用黏合剂的品种不同 ， 对药物的崩解和溶出的影响也不一样 。 
    \begin{blankblock}{}
        选用三种不同黏合剂 （ 10\% 阿拉伯胶浆 、 5\%淀粉浆 、 2\%聚维酮乙醇溶液 ） 制备氢氯噻嗪片 ，其崩解度和溶出速率都以淀粉浆制粒的片剂最快 ， 其余依次是 10\% 阿拉伯胶浆和 2\% 聚维酮乙醇溶液 。
    \end{blankblock}

\end{frame}


\begin{frame}
    \frametitle{崩解剂}
    片剂中加入崩解剂的主要目的是消除因黏合剂或由于加压而形成的结合力而使片剂崩解 。 崩解剂的品种和用量会对药物的溶出产生影响 。常用的崩解剂有：干淀粉、羧甲基淀粉钠、低取代羟丙基纤维素、交联羧甲基纤维素钠、交联聚维酮、泡腾崩解剂等。

        \begin{figure}[htpb]
            \centering
            \includegraphics[width=.3\linewidth]{崩解}
            \caption{不同剂量的崩解剂 （ 淀粉 ） 与水杨酸钠片溶出度的关系}
        \end{figure}
\end{frame}


\begin{frame}
    \frametitle{润滑剂}
    片剂压片时为了能顺利加料和切片，并减少黏冲及颗粒与颗粒间、药片与模孔间的摩擦力而加入的辅料称润滑剂。在制备胶囊剂时，粉末的流动性对装填效果很大。

    常用的润滑剂多为疏水性或水不溶性物质。疏水性润滑剂包裹在颗粒表面 ， 使水分不易渗入 ， 影响片剂的崩解与溶出 ； 而亲水性润滑剂能够促进药物与胃肠液的接触 ， 使集结的颗粒分散到胃肠液中 ， 则能使药物溶出量大幅度增加 。 正确地把握润滑剂对药物溶出的影响，已成为胶囊剂处方设计的重要课题。
    \begin{blankblock}{}
        硬脂酸镁与滑石粉为常用的润滑剂 ， 前者具有疏水性 ， 后者为水不溶性物质但具有亲水性 ， 度米芬咽喉片用硬脂酸镁作润滑剂时 ， 最低抑菌浓度为 1 ： 3860 ， 改用滑石粉后 ，最低抑菌浓度为 1：100000 ， 提高约 25 倍 。
    \end{blankblock}
\end{frame}


% \begin{frame}
%     \frametitle{增稠剂}
%     溶液剂及混悬剂等液体制剂中常加入一些增稠剂以改善制剂的流变学性质 、 物理稳定性等 。 药物的溶出度和扩散速度与黏度呈反比 ， 溶液剂黏度增加可能减缓胃排空 ， 或减缓药物分子到达吸收表面的扩散速度等 ； 混悬液黏度增加可能减缓药物溶解 。 因此制剂黏度增加往往会影响药物的吸收 。 

%     \begin{blankblock}{}
%         大鼠分别灌胃给予水杨酸钠溶液和含 2\% 甲基纤维素的水杨酸钠溶液后 ， 前者在血浆中和脑中出现的速度较后者快 ； 
        
%         又如增加苯巴比妥钠溶液中的蔗糖浓度时 ， 大鼠服用后可延长麻醉的诱导期 。
%     \end{blankblock}
% \end{frame}


\begin{frame}
    \frametitle{表面活性剂}
    表面活性剂在药剂中使用较为广泛，往往会对药物吸收产生影响。可促进药物的吸收，也可延缓药物的吸收。

    低浓度时：
    \begin{enumerate}
      \item 降低表面张力 ，因而增加疏水性药物的湿润性，使固体药物与胃肠液的接触角变小，提高有效表面积 ，因而增加药物的溶出和吸收 。 
      \item 溶解上皮细胞膜脂质，提高上皮细胞膜的渗透性 ，使吸收增加 。
    \end{enumerate}

    高浓度时$>$CMC：
    \begin{enumerate}
        \item 表面活性剂能形成胶束 ，药物包裹在胶团中，减少吸收 。
        \item 使细胞膜崩解或溶解 ，使部分膜蛋白变性或以薄层包在细胞膜周围 ， 从而干扰吸收 。
      \end{enumerate}
\end{frame}


\begin{frame}
    \frametitle{表面活性剂}
    \hspace{-2em}
    \begin{figure}[htpb]
        \centering
        \setcounter{subfigure}{0}
        \subfloat[低浓度聚山梨80对非那西丁吸收的影响]{\includegraphics[height=5.2cm]{表面活性剂}}\ 
        \subfloat[高浓度聚山梨80对大鼠小肠中的水杨酰胺吸收的影响]{\includegraphics[height=5.2cm]{高浓度}}
    \end{figure}
\end{frame}


\begin{frame}
    \frametitle{\thesection.\thesubsection.3 药物间及药物与辅料的相互作用}
    随着药品品种的增多及临床药学的普及深入，临床治疗中合并用药已很普遍。合并用药的目的在于提高药物的治疗效果，减少或消除药物的不良反应。

    \begin{itemize}
      \item 络合作用
      \item 吸附作用
      \item 包合作用
    \end{itemize}
\end{frame}


\begin{frame}
    \frametitle{络合作用}
    药物在制剂中可能与辅料发生相互作用 ， 如络合作用 、 吸附作用以及形成胶束等 ， 都能使药物在吸收部位的浓度降低 。 药物的络合物通过氢键结合 ， 如果该络合物在体液中能够溶解 ， 则说明两个部分间的作用是可逆的 。 药物络合物的性质 ， 如溶解度 、分子大小 、 扩散性以及油 / 水分配系数 ， 可能与原来的药物有很大差别 。 药物络合物中被络合的药物是以不能被吸收的形式存在的 ， 使药物的有效浓度降低 。 
    \[
        \text{药物} + \text{络合剂} \rightleftharpoons  \text{药物络合物}
    \]
    络合作用的程度用稳定常数$K_\mathrm{s}$表示 ，络合作用对吸收的影响取决于$K_\mathrm{s}$的大小 。
    \[
        K_\mathrm{s}=\frac{[\text{ 药物络合物 }]}{[\text{ 药物 }]\times[\text{ 络合物 }]}
    \]
\end{frame}


\begin{frame}
    \frametitle{络合作用}
    \begin{blankblock}{}
        制剂中广泛使用的高分子聚合物如树胶 、 纤维素衍生物 、 多元醇类及非离子型表面活性剂等 ， 与药物间的络合作用一般是可逆反应 ， 故对药物的吸收影响较小 ， 但也有例外 ，如苯丙胺与羧甲基纤维素形成难溶性络合物 ， 使其生物利用度大幅度下降 。 所以 ， 苯丙胺制剂不宜用纤维素类衍生物作为混悬剂或黏合剂 。 
    \end{blankblock}

    \begin{blankblock}{}
        苯巴比妥与 PEG4000 可形成一种溶解度很低的不吸收的络合物 ， 使含 PEG4000 的苯巴比妥片剂溶出速率大为减小 。
    \end{blankblock}

    \begin{blankblock}{}
        含二价或三价金属离子(Ca$^{2+}$、 Mg$^{2+}$ 、 Fe$^{3+}$、 Al$^{3+}$等) 化合物与四环素类抗生素或喹诺酮类抗生素同时服用，在胃肠道可形成难溶络合物，使抗生素吸收受阻 。 
    \end{blankblock}
\end{frame}


\begin{frame}
    \frametitle{吸附作用}
    \begin{enumerate}
      \item 物理吸附
      
      物理吸附指从溶液中将药物分子除去并转移到 “ 活性 ” 固体表面 ， 溶液中药物与被吸附药物间常存在平衡关系 。 
      \item 化学吸附
      
      如果吸附是不可逆的 ，表明药物与 “ 活性 ” 固体表面存在很强的键合作用 ， 则为化学吸附 ， 化学吸附无疑对药物吸收产生显著影响 。 
    \end{enumerate}
    
    \begin{blankblock}{}
        活性炭能够    吸附多种药物 ， 如抗生素 、 激素类 、 生物碱类等 。 
        
        药物与白陶土制剂同时服用 ， 则药物的吸收会减少 ， 相应的血药浓度也会降低 。
    \end{blankblock}
\end{frame}


% \begin{frame}
%     \frametitle{固体分散作用}
%     固体分散作用可加快药物的溶出 ， 也能延缓药物的释放 ， 依赖于所使用载体材料的性质 。
%     \begin{itemize}
%       \item 如果药物以分子状态 、 胶体状态或微品状态分散于水溶性载体中 ， 可构成一种高度分散体系 ， 从而增加难溶性药物的溶出速率和吸收速率 。
%       \item 倘若用疏水性 、 肠溶性或脂质类材料为载体制备固体分散体 ， 由于载体材料能形成可容纳药物的网状骨架结构 ， 被分散在骨架内的药物分子或微晶必须通过网状结构慢速扩散而溶出 ， 使整个释放过程减慢 ， 药物的吸收受释放过程控制而缓慢 。
%     \end{itemize}
% \end{frame}


\begin{frame}
    \frametitle{包合作用}
    将药物分子包嵌于另一种物质分子的空穴结构内的制剂技术称为包合作用 。 包合物的形成可视为药物与包合材料产生了相互作用的结果 。包合物由主分子和客分子两部分组成 。主分子为具有一定空穴结构的药用材料 ， 小分子药物作为客分子被包合在主分子内 ，形成分子囊 。
    \begin{blankblock}{}
        药剂学中常用的包合材料为$\beta$环糊精($\beta$-CYD)，它外端具有良好亲水性，内部具有疏水性，从而具有类似表面活性剂的性质，包合药物后形成水溶性较大的包合物。脂溶性药物的疏水键与$\beta$-CYD空洞中疏水键相互作用 ， 极性药物分子与$\beta$-CYD的羟基形成氢键结合 。 药物被水溶性主分子包合后 ， 通常药物的溶解度和溶出速率得到改善 ， 使药物的吸收增加 。
    \end{blankblock}
    
\end{frame}


\begin{frame}
    \frametitle{\thesection.\thesubsection.4 制剂工艺与药物吸收}
    \begin{enumerate}
      \item 混合与制粒
      \item 压片
      \item 包衣
    \end{enumerate}
\end{frame}


\begin{frame}
    \frametitle{混合与制粒}
    混合方法不同易引起药物溶出速率的差异 ， 尤其是对于小剂量的药物影响更明显 。 
    
    粉体性质 （ 如粒子的粒径 、 形态 、 密度等 ） 、 混合方式 、 混合时间 、 操作条件及设备等都会影响混合效果 。 如用溶媒分散法将剂量小的药物配成溶液再与辅料混合 ， 比将药物直接与辅料混合分散均匀度好得多 ， 亦有利于药物的溶出 。 

    即使是同样的处方 ， 制粒方法不同 ， 不仅所得颗粒的形状 、 大小 、 密度和强度不同 ， 而且其崩解性 、 溶解性也可能有很大差别 ， 药物疗效会受到影响 。
    \begin{itemize}
      \item 干法制粒
      \item 湿法制粒
    \end{itemize}
\end{frame}


\begin{frame}
    \frametitle{压片}
    压片是在压力下将颗粒状或粉末状药物压实的过程 。 压力的大小影响片剂的孔隙率 ， 进而影响片剂的崩解与药物的溶出 。 
    \begin{itemize}
      \item 一般情况下 ， 压力增大 ， 片剂的孔隙率减小、硬度变大、比表面积变小，崩解时间延长 ， 溶出速率变慢 。
      \item 压力与溶出速率的关系还与原辅料有关。塑性较强的物料受压时易产生塑性变形 ， 可
      压性好 ， 压制的片剂硬度比较大 。
    \end{itemize}

    \vspace{-2pt}
    \begin{blankblock}{}
        用磷酸氢钙压片时 ， 压力在一定的范围内 ， 片剂的比表面积随压力增大而逐渐增大 ， 溶出速率加快 ；而用微晶纤维素压片时 ， 压力增大 ， 溶出速率减小 。 因为微晶纤维素受压时粒子结合即发生塑性形变 ， 所以压力增大 ，孔隙率及比表面积减小 ， 溶出速率也就降低 。
    \end{blankblock}
     
\end{frame}


\begin{frame}
    \frametitle{包衣}
    包衣材料性质 、 包衣液组成 、 包衣层厚度等与包衣相关的因素可影响包衣制剂的溶出行为 ， 因而影响药物的吸收的快慢及血药浓度的高低 。
    \begin{figure}[htpb]
        \centering
        \includegraphics[width=.44\linewidth]{包衣}
        \caption{包衣对阿司匹林片溶出速率的影响}
    \end{figure}
\end{frame}


\section{口服药物吸收的评价方法}

\begin{frame}
    \frametitle{口服药物吸收的评价方法}
    口服药物吸收评价主要通过药物渗透性进行。药物渗透性的测定可采用人体试验方法或其他能预测药物体内吸收程度的非人体试验方法。
    \begin{enumerate}
      \item 体内法（in vivo）
      \item 在体法（in situ）
      \item 体外法（in vitro）
    \end{enumerate}
\end{frame}


\begin{frame}
    \frametitle{体内法}
    体内法是以整体动物机体为研究对象，进行药动学研究。

    口服给药后，在不同时间点采集生物样本如血液、尿液，再通过一定方法测定其药物浓度，计算达峰时间（$T _\mathrm{max}$）、达峰浓度（$C_\mathrm{max}$）、药时曲线下面积（$AUC$）等药动学参数，进而评价药物吸收的程度和速度。

    \begin{itemize}
      \item 给药方法一般是大动物口服给药，小动物灌胃给药。
      \item 此法采用整体动物进行实验，能真实地反映药物口服后的体内吸收情况，但是它综合了物化、生理、剂型等众多因素的结果，不能特异性地反映肠道对药物的吸收作用。
      \item 由于实验周期长、操作相对繁琐、影响因素较多、动物个体差异较大等原因，故体内法很少应用于药物吸收机制的研究，一般用于研究药物体内药动学特征。
    \end{itemize}

\end{frame}


\begin{frame}
    \frametitle{在体法}      
    在体回肠灌流法：在动物麻醉状态下，不切断血管和神经，将待研究的肠段两端结扎，肠管插管，洗净肠内容物，构成循环回路，使药物在肠腔内循环。不同时间测定灌流液中药物浓度变化，可以获得药物在肠内吸收情况。本法能避免胃容物和消化道固有生理活动对结果影响。
      \begin{itemize}
        \item 在体法是建立在整体动物水平上的实验，与体内法不同的是其将干扰因素大大减少，同时保证了肠道神经与内分泌调节的完整性，也保证了血液与淋巴液供应，使得待考察组织的生物活性大大提升，能够较为准确地反映药物在肠道的真实吸收情况，常用于研究药物的渗透和吸收动力学。
        \item 在体法由于是建立在整体水平之上，因此个体差异较大，对实验动物的数量有一定要求，以保证最小的数据变异。
      \end{itemize}

\end{frame}


\begin{frame}
    \frametitle{体外法}
    \begin{enumerate}
        \item 肠外翻囊法
        
        动物麻醉后， 取出一段肠管，用pH 7.40 Krebs-Henseleit (K-H)生理液冲洗，将肠管套于玻璃棒上外翻， 一端固于聚乙烯管上， 另一端结扎，内注入一定量的pH 7.40 K-H生理溶液，置于含有药物的pH 7.40 K-H生理溶液中，37℃恒温， 在5\%O$_2$气流下，定时从内外两侧取样，测定药物浓度变化。该方法可根据需要研究不同肠段的药物吸收或分泌特性及其影响因素。
        \item 细胞培养模型法
      \end{enumerate}
      \begin{itemize}
        \item 操作简单、试验周期短，可以用于药物早期高通量筛选。
        \item 肠道的不同节段以及缺乏血液、淋巴液的供应对细胞旁路通道和酶活性造成的影响均会影响实验结果，其组织活性与其在体内环境中会存在较大的差异，不能反映药物在肠道环境中的真实吸收情况。
      \end{itemize}
\end{frame}


\begin{frame}
    \frametitle{Caco-2细胞模型}
    细胞模型是体外法评价药物吸收的重要手段，最常用的是Caco-2细胞模型。

    Caco-2细胞模型(human colon adenocarcinoma
cell line ) 最早是由 美国 堪萨斯大学 药学 院
Ronald T. Borchardt 博士的实验室以及英国 的
Ciba-Geigy实验室在二十世纪八十年代后期建立的。
是第一个被国外制药企业以及实验室广泛应用的肠
道吸收的细胞模型，也是迄今为止的最好的肠道转
运模型和上皮转运模型。

Caco-2细胞系来源于人的直肠癌， 其结构和生
化作用类似于人小肠上皮细胞， 含有与小肠刷状缘
上皮相关的酶系， 与正常的成熟小肠上皮细胞在体
外培育过程中出现逆向分化(反分化) 不同， Caco-2
细胞在传统的细胞培养条件下， 生长在多孔的可渗
透的聚碳酸酯膜上可达到融合并自发分化为肠上皮
细胞， 形成连续的单层(monolayer) 。培育约3星期后， 其电阻值达到400 Ω （有的
要求达到600Ω ） ， 它们形成致密的细胞单层，
结构和生化作用类似于小肠上皮细胞， 3-5天
内适合进行转运试验。
\end{frame}


% \begin{frame}
%     \frametitle{Caco-2细胞模型}
% Caco-2细胞单层膜模型构建成功后，在膜的两侧分为了A侧（肠腔侧）与B侧（基底侧）。

% \begin{figure}[htpb]
%     \centering
%     \includegraphics[width=.8\linewidth]{caco2}
% \end{figure}
% \end{frame}


\begin{frame}
    \frametitle{Caco-2细胞模型的优点}
    \begin{enumerate}
      \item Caco-2细胞来源是人结肠癌细胞，同源性好。该模型与细胞在肠中的吸收有良好的相关性、较高的重现性；
      \item 在细胞水平上研究药物在小肠粘膜中的吸收、转运和代谢
      \item 与动物实验相比， Caco-2细胞生命力强， 培养方法简单，更经济高效；
      \item 可同时对大批量药物进行快速筛选， 所需药物量少， 仅需要数毫克即可；
      \item 可以从不同方向进行摄取和转运实验，可用于区分肠腔内不同吸收途径的差别；
      \item 有较广泛的适用性（可用于原料药、制剂、大小分子的跨肠膜研究）
      \item 可以同时研究药物对粘膜的毒性
    \end{enumerate}
\end{frame}


\begin{frame}
    \frametitle{Caco-2细胞模型的局限}
    \begin{enumerate}
      \item 缺少粘蛋白产生细胞， 缺少肠壁粘液层；
      \item 具有较高的细胞通透性， 若喂养次数多， 某些酶载体或细胞内容物溢出；
      \item 缺少某些代谢酶；
      \item 细胞的培养和代系对于Caco-2细胞形态学和生理性质有影响， 引起细胞对药物转运产生差别；
      \item 由于培养条件差别， 使不同实验室出来结果缺乏可比性
    \end{enumerate}
\end{frame}


\section{药物在其它部位的吸收}


\begin{frame}
    \frametitle{在口腔粘膜中吸收}
    药物经口腔黏膜给药可发挥局部或全身治疗作用，局部作用的剂型多为溶液型或混悬型，如漱口剂 、气雾剂 、膜剂、软膏剂、口腔片剂等 。 欲产生全身作用常用舌下片 、黏附片、贴剂等 。产生全身作用的口腔用药主要部位是颊黏膜，其次是舌下黏膜 。

    \begin{figure}[htpb]
        \centering
        \includegraphics[width=.5\linewidth]{口腔给药}
    \end{figure}
\end{frame}

\begin{frame}
    \frametitle{在口腔粘膜中吸收}
    \hspace{-2em}
    \begin{minipage}{0.28\linewidth}
        \begin{figure}[htpb]
            \centering
            \includegraphics[width=\linewidth]{颊粘膜}
            \caption{颊粘膜结构示意图}
        \end{figure}
    \end{minipage}\hfill
    \begin{minipage}{0.68\linewidth}
         颊黏膜可分为四层结构 ： 
        \begin{enumerate}
          \item 上皮层：由角质形成细胞与非角质形成细胞组成 ， 为药物透过黏膜的主要屏障 。
          \item 基底层
          \item 固有层：致密的结缔组织
          \item 黏膜下层：疏松的结缔组织
        \end{enumerate}
        药物在口腔黏膜中渗透到达黏膜下层 ， 黏膜下层有大量的毛细血管汇总至颈内静脉而进入血液循环 ， 因此可避免肝脏首过效应 。
    \end{minipage}
\end{frame}

\begin{frame}
    \frametitle{在直肠粘膜中吸收}
    直肠的吸收面积不大，但血流量较为丰富，药物容易吸收。直肠给药可用于局部治疗或发挥全身作用 ， 常用的剂型是栓剂或灌肠剂 。 栓剂用于全身治疗有许多优点 ： 
    \begin{enumerate}
      \item 药物直肠吸收后 ， 大部分绕过肝脏进入大循环发挥全身作用 ， 降低了肝首过效应 ， 也相应减少药物对肝脏的毒副作用； 
      \item 避免胃肠 pH 和酶的影响和破坏 ， 避免药物对胃肠功能干扰 ， 对胃有刺激的药物可采用直肠给药 ；
      \item 作用的时间一般比片剂长 ， 通常给药为 1 、 2 次 / 日；
      \item 适于不愿或不能吞服药物的患者 ， 对于婴幼儿及神志障碍的患者 ， 使用栓剂较口服或注射给药更容易 、 更安全 ；
      \item 可作为多肽蛋白质类药物的吸收部位 。
    \end{enumerate}
\end{frame}

\begin{frame}
    \frametitle{在直肠粘膜中吸收}
    直肠给药制剂首先需要药物从基质中释放，释放的药物穿过上皮细胞进入血管被吸收。
    \begin{figure}[htpb]
        \centering
        \subfloat[药物从基质中的释放与吸收过程]{\includegraphics[height=5cm]{直肠给药}}\quad
        \subfloat[直肠与肛门部位的血管分布]{\includegraphics[height=5cm]{直肠血管分布}}
    \end{figure}
\end{frame}

\begin{frame}
    \frametitle{透皮肤吸收}
    皮肤给药可以用于局部皮肤病的治疗，也可以经皮肤吸收后治疗全身性疾病。对于皮肤病，由于病灶部位的深浅不同 ，某些药物需要透过角质层以后才能起效；而对于全身性疾病 ， 药物必须通过角质层 ，被皮下毛细血管吸收进入血液循环后才能起效。经皮吸收是指药物从应用于皮肤上的制剂中释放与穿透皮肤进入体循环的过程 。

    \hspace{-2em}
    \begin{minipage}{0.38\linewidth}
      \begin{figure}[htpb]
        \centering
        \includegraphics[width=\linewidth]{皮肤结构}
      \end{figure}
    \end{minipage}\hfill
    \begin{minipage}{0.58\linewidth}
      \begin{itemize}
        \item 角质层是经皮吸收药物的主要屏障。
        \item 当皮肤角质层受损时，药物的通透性显著增加。
        \item 不同部位的角质层厚度不同，不同的药物在不同的部位吸收程度是不同的。
        \item 皮肤有较大的表面积, 皮肤吸收往往也是毒物（如农药）等吸收造成中毒的重要途径之一。 
      \end{itemize}
    \end{minipage}
\end{frame}


\begin{frame}
    \frametitle{在肺部中吸收}
    肺部给药又称吸入给药， 主要是通过口腔吸入 ， 经过咽喉进入呼吸道 ， 到达呼吸道深处或肺部 ， 起到局部作用或吸收后产生全身治疗作用 。 

    \begin{itemize}
      \item 气体或挥发性药物可以在肺上皮细胞或气管粘膜吸收
      \item 用于肺部给药的剂型包括气雾剂 、 喷雾剂和粉雾剂 
      \item 肺吸收也往往是毒品和环境毒物重要的吸收途径
      \item 优点是药物吸收快，可避免肝肠分解
      \item 缺点是药物的剂量难以控制，药物也会对肺上皮细胞产生刺激
    \end{itemize}
\end{frame}


\begin{frame}
  \frametitle{在肺部中吸收}
  \hspace{-2em}
  \begin{minipage}{0.4\linewidth}
    \begin{figure}[htpb]
      \centering
      \includegraphics[width=\linewidth]{呼吸系统}
      \caption{人体呼吸器官示意图}
    \end{figure}
  \end{minipage}\hfill
  \begin{minipage}{0.58\linewidth}
    \begin{itemize}
      \item 气体进入口腔后经咽喉部 、 气管 、 左右支气管进入左右肺叶 ，进入肺叶后支气管进一步细分为更细的支气管 ， 而后进入终末细支气管 。
      \item 从终末细支气管后继续分支 ， 每次分支为两条导管 ， 最后至肺泡管并与肺泡相连 。
      \item 正常人的肺部表面积为 \SI{100}{m^2}
      \item 肺泡是血液与气体进行交换的部位
      \item 巨大的肺泡表面积、丰富的毛细血管和极小的转运距离，决定了肺部给药迅速吸收，且吸收后的药物直接进入血液循环，不受肝首过效应影响。
    \end{itemize}
  \end{minipage}
\end{frame}


\begin{frame}
    \frametitle{在眼部吸收}
    眼部给药主要用于眼局部疾病的治疗 ， 如抗眼部细菌性或病毒性感染 、 降低眼压 、 缩瞳或扩瞳等 。 眼部给药后药物能够到达眼内病灶部位 ， 发挥疾病的治疗作用 。眼睛是人的重要器官 ， 又非常敏感 ， 因此一般不会作为全身治疗作用的给药途径 。
    \hspace{-2em}
    \begin{minipage}{0.38\linewidth}
      \begin{figure}[htpb]
        \centering
        \includegraphics[width=\linewidth]{眼结构}
      \end{figure}
    \end{minipage}\hfill
    \begin{minipage}{0.58\linewidth}
      药物在眼部吸收分为角膜吸收和非角膜吸收两个途径。
      \begin{enumerate}
        \item 角膜吸收一般是眼局部用药的有效吸收途径，角膜外层上皮细胞是药物的主要转运屏障。
        \item 非角膜吸收是药物进入血液循环的主要途径。
      \end{enumerate}
    \end{minipage}
\end{frame}


\begin{frame}
    \frametitle{在眼部吸收}
    \hspace{-2em}
    \begin{minipage}{0.48\linewidth}
      药物在眼部吸收有以下特点：
      \begin{enumerate}
        \item 简单、经济
        \item 可以避免肝首过效应
        \item 对免疫反应不敏感
      \end{enumerate}

      ~~
    \end{minipage}\hfill
    \begin{minipage}{0.48\linewidth}
      眼部给药存在以下缺陷：
      \begin{enumerate}
        \item 眼部刺激性
        \item 剂量损失
        \item 药物在眼部的停留时间问题
        \item 一般病人难以接受眼部给药
      \end{enumerate}
    \end{minipage}
\end{frame}


\begin{frame}
    \frametitle{在鼻腔粘膜中吸收}
    鼻腔给药不仅适用于鼻腔局部疾病的治疗 ， 也是全身疾病治疗的给药途径之一 。 鼻腔给药的药物吸收是药物透过鼻黏膜向循环系统的转运过程。 鼻腔给药的主要优点有 ：

    \begin{enumerate}
      \item 鼻黏膜薄 、 有效表面积大 、 渗透性高 、 血管丰富 ， 有利于吸收 ， 吸收程度和速度有时可与静脉注射相当 ；
      \item 药物直接进入血液循环，可避开药物在胃肠液中的降解和肝首过作用 ；
      \item 能够增加一些药物向脑内传递 ， 有利于脑部疾病的治疗；
      \item 鼻腔内给药方便易行
    \end{enumerate}

    但鼻腔给药也存在不足 ， 如单次用剂量有限 ， 某些药物生物利用度相对较低 ， 吸收剂量不够准确 ， 药物与吸收部位的接触时间相对较短等 。
\end{frame}


\begin{frame}
    \frametitle{在鼻腔粘膜中吸收}
    口服给药个体差异大 、 生物利用度低的药物以及口服易破坏或不吸收 、 只能注射给药的药物 ， 可考虑鼻黏膜给药 。 如在胃肠道中难吸收的磺苄西林 、 头孢唑林可经鼻黏膜吸收 ， 维生素
    \begin{blankblock}{}
      在胃肠道中难吸收的磺苄西林 、 头孢唑林可经鼻黏膜吸收 ， 维生素B12 的鼻用凝胶剂比同剂量口服片剂血药浓度高 $8\sim 10$ 倍 。 药物可制成溶液剂滴入鼻腔 ， 也可以气雾剂给药 。 

      用来治疗头痛和偏头痛的酒石酸布托啡诺鼻喷剂可达到与静脉注射相当的血药浓度水平。
    \end{blankblock}
\end{frame}


\begin{frame}
  \frametitle{在肌肉、皮下吸收}
  \hspace{-2em}
  \begin{minipage}{0.3\linewidth}
    \begin{figure}[htpb]
      \centering
      \includegraphics[width=\linewidth]{注射给药}
      \caption{注射给药示意图}
    \end{figure}
  \end{minipage}\hfill
  \begin{minipage}{0.68\linewidth}
    \begin{enumerate}
      \item 静脉注射药物直接注入血液循环 ， 不存在吸收过程 。静脉注射的生物利用度一般被认为是 100\% 。
      \item 肌内注射存在吸收过程 ， 药物先经注射部位的结缔组织扩散 ， 再经毛细血管吸收进入血液循环 ， 所以肌内注射药物的起效比静脉注射稍慢 ， 但比静脉注射简便安全 ， 比皮下注射刺激性小 ， 应用较广 。肌肉组织内的血管十分丰富 ，一般认为肌内注射给药的药物吸收程度与静脉注射相当 。
    \end{enumerate}
  \end{minipage}

  \begin{enumerate}
    \item[3.] 皮下注射是将药物注射到疏松的皮下组织中 。皮下组织血管较少 ， 血流速度比肌肉组织慢 。 因此 ， 皮下注射后药物吸收较肌内注射慢 ， 有时甚至比口服吸收还慢。皮下植入给药往往可维持数周或数月的疗效。 需延长作用时间的药物可采用皮下注射 ， 如治疗糖尿病的胰岛素等 。
  \end{enumerate}
\end{frame}


\begin{frame}[t]
    \frametitle{}
    \begin{question}
        不存在吸收过程的给药途径是
        \tcblower

        A. 口服给药

        B. 肌内注射

		    C. 肺部给药

        \only<1>{
            D. 静脉注射
        }
        \only<2>{
            \highlight{D. 静脉注射}
        }
    \end{question}
\end{frame}

\end{document}